\chapter{MGTP $BF~NO8@8l;EMM(B}
$B$3$N>O$G$O!"(BMGTP $B$NF~NO8@8l$NI8=`E*%7%s%?%/%9$K$D$$$F@bL@$7!"(B
$B5-=RNc$r<($9!#(B

%%----------------------------------------------------------------
\section{MGTP $BF~NO@a%7%s%?%/%9(B}

MGTP $B$NF~NO@a$O!"Bh(B1$B>O$N5!G=@bL@$N$H$3$m$G$b@bL@$7$?$h$&$K!"(B
$B@5@a(B(Positive clause)$B!"Ii@a(B(Negative clause)$B!":.9g@a(B(Neutral clause)
$B$N(B3$B<oN`$,$"$k$,!"4pK\E*$K$9$Y$F$N@a$O4^0U(B(implicatoin) \verb+-->+ $B$K(B
$B$h$C$F9=@.$5$l$F$$$k!#(B
$B4^0U$N:8JU$rA07oIt(B(Antecedent)$B!"8e7oIt$r(B(Consequent) $B$H8F$V!#(B

$B@5@a$N>l9g$O!"A07oIt$r(B \verb+true+ $B$K!"$^$?Ii@a$N>l9g$O!"(B
$B8e7oIt$r(B \verb+false+ $B$H$7$F5-=R$9$k!#(B
$BCM0h8BDj$N>r7o$rK~B-$7$J$$$h$&$J@a$K$D$$$F$O!"$"$i$+$8$a(B
\verb+dom+ $B$J$I$N=R8l$rMQ$$$k$3$H$K$h$C$F>r7o$r(B
$BK~B-$9$k$h$&$J@a$KJQ49$9$k$3$H$,2DG=$G$"$k!#(B

$B$^$?!"A07oIt!"8e7oIt$K$*$$$F$O!"3FBP>]8@8l$rD>@\8F$S=P$9(B
$B$3$H$b5v$5$l$F$*$j!"$3$l$K$h$j@0?t4V$NHf3S1i;;$J$I$,(B
$B2DG=$K$J$k!#(B

$B0J2<$O!"(BMGTP $BF~NO@a$N%7%s%?%/%9$K$D$$$F(B BNF $B$G5-=R$7$?$b$N$G$"$k!#(B

\vspace{1em}
\begin{tabular}{ll}
 MGTP-Clause ::= & Positive-Clause $|$ Negative-Clause $|$ Neutral-Clause\\
 Positive-Clause ::= & ' \verb+true+' ' \verb+-->+' Consequent '.'\\
 Negative-Clause ::= & Antecedent  ' \verb+-->+' ' \verb+false+' '.'\\
 Neutral-Clause ::= & Antecedent ' \verb+-->+' Consequent '.'\\
 Antecedent ::= & Antecedent-Atom { ',' Antecedent-Atom }*\\
 Antecedent-Atom ::= & MGTP-Atom $|$ ' \verb+{{+' [ Guards ] ' \verb+}}+'\\
 Consequent ::= & [ '\verb+{{+' [ Procedures ] '\verb+}}+' ' \verb+->+' ] Consequent-Group\\
 Consequent-Group ::= & Consequent-Unit { ',' Consequent-Unit }*\\
 Consequent-Unit ::= & Atom-Group $|$ Disjunctive-Atom-Group \\
 Disjunctive-Atom-Group ::= & Atom-Group { ' \verb+;+' Atom-Group }*\\
 Atom-Group ::= & MGTP-Atom { ',' MGTP-Atom }*\\
 MGTP-Atom ::= & Term \\
 Guards ::= & Call { ',' Call }*\\
 Procedures ::= & Call { ',' Call }*\\
\end{tabular}

\vspace{1em}
$B$3$3$G!"(BTerm $B$O!"(BProlog $B$d(B KL1 $B$N9`$rI=$9!#(B
$B$9$J$o$A=R8l$KAjEv$9$k!#JQ?t$NI=8=$b(B Prolog $B$d(B KL1 $B$N$=$l$K=`$:$k!#(B
$B$^$?!"(BCall $B$O3:Ev%W%m%0%i%`$N8F$S=P$7J8$r0UL#$7!"(B
$BNc$($P!"(BProlog $B$G$O(B \verb+{{X>Y}}+ $B$d(B \verb'{{X1:=X+1,X1>Y}}' $B$J$I$N(B
$B$h$&$J5-=R$,2DG=$G$"$k!#(B

$B$J$*!":Y$+$$%7%s%?%/%9$NAj0c$K$D$$$F$O3F%P!<%8%g%s$N%^%K%e%"%k$r;2>H(B
$B$5$l$?$$!#(B

%%----------------------------------------------------------------
\section{$BNcBj(B}

MGTP $BF~NO@a$NNc$r?^(B\ref{MGTP-input-ex}$B$K<($9!#(B
$B$3$NLdBj$O(B TPTP problem library\cite{SSY93} $B$N(B
MSC006-1($B?^(B \ref{MSC006-1}) $B$r(B MGTP $BF~NO@a$K=q$-D>$7$?$b$N$G$"$k!#(B

\begin{figure}
\small
\begin{verbatim}
%%  Problems/MSC/MSC006-1.mg  Ground-MGTP/pl input-clauses
%%  infinite_domain(no).
%%  max_number_of_doms_in_a_clause(2).
%%  total_number_of_doms(2).
%%  number_of_NonHorn_clauses(1).
%%  max_number_of_consequents_in_a_clause(2).
%%  number_of_negated_literals(0).

true-->dom(a),dom(b),dom(c),dom(d).
dom(A),dom(B)-->p(A,B);q(A,B).
p(A,B),p(B,C)-->p(A,C).
q(A,B),q(B,C)-->q(A,C).
q(A,B)-->q(B,A).
p(a,b)-->false.
q(c,d)-->false.

%%  Problems/MSC/MSC006-1.mg  EOF
\end{verbatim}
\caption{MGTP $BF~NO@a$NNc(B}
\label{MGTP-input-ex}
\end{figure}

\begin{figure}
\small
\begin{verbatim}
%--------------------------------------------------------------------------
% File     : MSC006=NonObv-1 : Released v0.0.0, Updated v0.11.5.
% Domain   : Miscellaneous
% Problem  : A ``non-obvious'' problem
% English  : Suppose there are two relations, P and Q. P is
%            transitive, and Q is both transitive and symmetric.
%            Suppose further the ``squareness'' of P and Q: any two
%            things are related either in the P manner or the Q
%            manner. Prove that either P is total or Q is total.
% Refs     : Pelletier F.J., and Rudnicki P. (1986), Non-Obviousness,
%            In Wos L. (Ed.), Association for Automated Reasoning
%            Newsletter (6), Association for Automated Reasoning,
%            Argonne, Il, 4-5.
% Source   : [Pelletier & Rudnicki, 1986]
% Names    : nonob.lop [SETHEO]
% Syntax   : Number of clauses           :    6 (   1 non-Horn)(   2 units)
%            Number of literals          :   12 (   0 equality)
%            Number of predicate symbols :    2 (   0 propositions)
%            Number of function symbols  :    4 (   4 constants)
%            Number of variables         :   10 (   0 singletons)
%            Maximal clause size         :    3
%            Maximal term depth          :    1
%--------------------------------------------------------------------------
input_clause(p_transitivity,hypothesis,
    [--p(X,Y),
     --p(Y,Z),
     ++p(X,Z)]).
input_clause(q_transitivity,hypothesis,
    [--q(X,Y),
     --q(Y,Z),
     ++q(X,Z)]).
input_clause(q_symmetry,hypothesis,
    [--q(X,Y),
     ++q(Y,X)]).
input_clause(all_related,hypothesis,
    [++p(X,Y),
     ++q(X,Y)]).
input_clause(p_is_not_total,hypothesis,
    [--p(a,b)]).
input_clause(prove_q_is_total,theorem,
    [--q(c,d)]).
%--------------------------------------------------------------------------
\end{verbatim}
\caption{MSC006-1(TPTP problem library)}
\label{MSC006-1}
\end{figure}

\normalsize

$B$^$?!"(B8 $B%/%$!<%sLdBj$r(B MGTP $B$G5-=R$7$?Nc$r(B
$B?^(B \ref{queen} $B$K<($9!#(B
$B=R8l(B \verb+p(M,N)+ $B$O(B \verb+M+ $B9T(B \verb+M+ $BNs$K(B
$B%/%$!<%s$,CV$+$l$k$3$H$rI=$7$F$$$k!#(B
$B$O$8$a(B$B$K@5@a$H$7$F!"3F9T$K$*$1$kCM$N8uJd$,A*8@$N7A<0$G(B
$B5-=R$5$l$F$$$k!#(B
$B$3$3$NNc$G$O!"(B\verb+{{ }}+ $B$G0O$^$l$?%W%m%0%i%`8F$S=P$7$,MQ$$$i$l$F$*$j!"(B
$B$3$l$K$h$j!"F1$89T$dF1$8Ns!"$"$k$$$OBP3Q@~>e$K%/%$!<%s$,CV$+$l$?(B
$B$3$H$rH=Dj$7!"%b%G%k8uJd$r4~5Q$9$k$h$&$K$J$C$F$$$k!#(B

\begin{figure}
\small
\begin{verbatim}
true --> p(1,1);p(1,2);p(1,3);p(1,4);p(1,5);p(1,6);p(1,7);p(1,8).
true --> p(2,1);p(2,2);p(2,3);p(2,4);p(2,5);p(2,6);p(2,7);p(2,8).
true --> p(3,1);p(3,2);p(3,3);p(3,4);p(3,5);p(3,6);p(3,7);p(3,8).
true --> p(4,1);p(4,2);p(4,3);p(4,4);p(4,5);p(4,6);p(4,7);p(4,8).
true --> p(5,1);p(5,2);p(5,3);p(5,4);p(5,5);p(5,6);p(5,7);p(5,8).
true --> p(6,1);p(6,2);p(6,3);p(6,4);p(6,5);p(6,6);p(6,7);p(6,8).
true --> p(7,1);p(7,2);p(7,3);p(7,4);p(7,5);p(7,6);p(7,7);p(7,8).
true --> p(8,1);p(8,2);p(8,3);p(8,4);p(8,5);p(8,6);p(8,8);p(8,8).

p(N,M),p(N1,M),{{N \= N1}} --> false.
p(N,M),p(N,M1),{{M \= M1}} --> false.
p(N,M),p(N1,M1),{{N \= N1,M \= M1,A:=N-N1,B:=M-M1,A=B}} --> false.
p(N,M),p(N1,M1),{{N \= N1,M \= M1,A:=N-N1,B:=M1-M,A=B}} --> false.
\end{verbatim}
\caption{8 $B%/%$!<%sLdBj$N(B MGTP $B$G$N5-=RNc(B}
\label{queen}
\end{figure}


\begin{thebibliography}{10}
\bibitem{SSY93}
Suttner, C., Sutcliffe, G. and Yemenis, T.:
%The TPTP Problem Library.
%TR12.11.93, 1993.
The TPTP Problem Library,
{\it Proc. 12th CADE,} pp. 252-266 (1993).
\end{thebibliography}
